% ID: MAT-TRI-EQG-016
% Titolo: Sistema goniometrico (1)
% Disciplina: Matematica
% Area: trigonometria
% Argomento: equazioni_goniometriche
% Classe: Quarta secondaria di secondo grado
% Difficolta: 3
% Tipo: problema_numerico
% Tag: matematica, trigonometria, equazioni_goniometriche
% Fonte: verifica_goniometria_base_anonimizzata
% Provenienza: materiale fornito dal docente; testo anonimizzato; soluzione redatta per Exergo

\begin{esercizio}
Risolvi il seguente sistema: 
\[
\begin{cases}
2\sin^2 x-\sin x-1=0,\\[2mm]
\cos x-\sin x=\dfrac{\sqrt{3}+1}{2}.
\end{cases}
\]
Salvo diversa indicazione, lavora in $\mathbb R$; angoli in radianti.
\end{esercizio}

\begin{soluzione}
Dalla prima: $(2\sin x+1)(\sin x-1)=0$.
Il caso $\sin x=1$ non verifica la seconda equazione.
Con $\sin x=-1/2$, la seconda impone $\cos x=\sqrt3/2$.
Quindi $x=11\pi/6+2k\pi$. In tutte le famiglie di soluzioni $k\in\mathbb{Z}$.
\end{soluzione}
